\chapter{State of the Art}
\label{state_of_art}

\section{BitTorrent Protocol}
\label{BitTorrent}

Traditional content distribution on the Internet relies on centralized architectures, typically using the HTTP protocol or Content Delivery Networks (CDNs). In this model, the entire bandwidth and availability cost falls upon the server responsible for hosting and delivering the content. Although CDNs can improve performance and scalability, the system still depends on a limited number of central nodes. If the server becomes unavailable, the content can no longer be retrieved.

To mitigate this limitation, peer-to-peer (P2P) architectures were introduced, allowing users to share bandwidth and storage resources directly. Among several P2P protocols developed over the years, \textbf{BitTorrent}, proposed by Bram Cohen in 2001, remains one of the most widely adopted and efficient solutions for large-scale content distribution. BitTorrent decentralizes the distribution process by dividing files into smaller fixed-size pieces (or \textit{chunks}) that can be shared among participants, known as \textit{peers}, within a swarm. Each peer downloads and uploads file pieces simultaneously, which increases redundancy, efficiency, and scalability.

\subsection{Torrent Metadata and File Structure}
\label{torrent_structure}

Before participating in a BitTorrent swarm, a client requires metadata describing the file it intends to download. This information is stored in a \textbf{.torrent} file, which contains:
\begin{itemize}
    \item The names and sizes of the files being shared.
    \item A list of SHA-1 hashes corresponding to each piece of the content.
    \item URLs of trackers, which help peers discover each other.
    \item Optional bootstrap nodes for Distributed Hash Table (DHT) discovery or Peer Exchange (PEX) connections.
\end{itemize}

This metadata allows peers to verify the integrity of each downloaded piece by computing its hash and comparing it with the expected value. Once all pieces are verified, the client can reconstruct the original file. Magnet links, commonly used today, encode the unique \textit{info hash} of the torrent and enable clients to obtain metadata through DHT without requiring a separate \textit{.torrent} file.

\subsection{Roles in the BitTorrent Swarm}
\label{bittorrent_roles}

BitTorrent defines several types of participants within a swarm, each playing a specific role in maintaining the network:

\begin{itemize}
  \item \textbf{Trackers:} Centralized or semi-centralized servers that facilitate peer discovery by maintaining a list of active peers in the swarm. They do not host any content; instead, they respond to peers’ \textit{announce} requests with a list of other peers participating in the same torrent. After the initial connection, peers communicate directly, while the tracker remains responsible for helping new peers join the swarm.

  \item \textbf{Downloaders (Leechers):} Peers that possess less than 100\% of the file. They download missing pieces from other peers while simultaneously uploading verified pieces they already own. The term “leecher” is sometimes used negatively to describe peers that download without contributing uploads, which harms overall swarm efficiency.

  \item \textbf{Seeders:} Peers that have obtained 100\% of the file and therefore only upload data to others. Seeders are crucial for the swarm’s survival, as at least one seeder must exist for other peers to complete their downloads.

  \item \textbf{Peers:} A general term for any node participating in the swarm, whether acting as a downloader or a seeder.
\end{itemize}

\subsection{Peer Discovery Mechanisms}
\label{peer_discovery}

In order to join a swarm, a peer must discover other participants. BitTorrent supports three main peer discovery mechanisms:

\begin{itemize}
  \item \textbf{Trackers:} The most traditional method, where the tracker provides a list of active peers associated with a given torrent. Communication with the tracker occurs through HTTP, HTTPS, or UDP.
  \item \textbf{Peer Exchange (PEX):} A decentralized mechanism allowing connected peers to exchange information about other peers in the swarm, reducing dependency on centralized trackers.
  \item \textbf{Distributed Hash Table (DHT):} A fully decentralized peer discovery system based on the Kademlia algorithm. Each peer acts as a node in the DHT, storing and retrieving peer information through key-value mappings. This mechanism enables trackerless operation and improves network resilience.
\end{itemize}

\subsection{Security and Privacy Issues in BitTorrent}
\label{issues_bittorrent}

Despite its efficiency and widespread use, BitTorrent suffers from several security and privacy vulnerabilities:

\begin{itemize}
  \item \textbf{Lack of Authentication:} Any peer can join a swarm without identity verification, allowing malicious nodes to distribute corrupted or fake content.
  \item \textbf{IP Address Exposure:} All peers publicly expose their IP addresses during communication, which enables traffic monitoring, de-anonymization, and legal tracking.
  \item \textbf{Sybil and Eclipse Attacks:} In DHT-based systems, attackers can create multiple fake identities or manipulate routing tables to control or isolate parts of the network.
  \item \textbf{Untrusted Client Implementations:} Since client software is open and decentralized, malicious or modified clients can misbehave, leak information, or disrupt swarm operations.
\end{itemize}

These limitations demonstrate the need for enhanced security and privacy mechanisms in decentralized file-sharing systems. In the following chapters, we explore how combining \textbf{Trusted Execution Environments (TEEs)}, the \textbf{QUIC} transport protocol, and \textbf{Onion Routing} can provide a more secure, anonymous, and trustworthy P2P file-sharing framework.
